% RESULTADO_RECUPERADO desde II REV09. Véase MAPA_PROCEDENCIA.json.
% Selección de líneas y cambios mecánicos explícitos; ninguna prueba nueva.
\section{Action, gravité et information thermique}
\label{v:ant:sec:gravity}
\subsection{Période dodécaphasique et réalisation circulaire de l'action}
La longueur de la section~\ref{g26:sec:rigidez-radial-es} fixe, dans
la même réalisation, l'horloge compatible
\[
 t_0=\frac{\pi_{\HMT}L}{54c_{\rm int}}.
\]
Sa représentation circulaire satisfait alors \(R_{108}=L\). Cette
substitution exprime chronologiquement la longueur construite et conserve
la référence longitudinale, l’action et la vitesse de cette réalisation.
Le calendrier déjà construit détermine $T_{108}=108t_0$.
Dans sa réalisation circulaire,
\[
\omega_{108}=\frac{2\pi}{108t_0},\qquad
R_{108}=\frac{c_{\rm int}}{\omega_{108}}
=\frac{54}{\pi}\ell_0,\qquad \ell_0=c_{\rm int}t_0.
\]
L’unité $\ell_0$ et les coordinatisations du temps et de la vitesse sont maintenues
explicites ; la coordinatisation SI ne sélectionne pas le cycle.
Pour chaque section orientée d’action $a\in\{\mathrm{pre},\mathrm{ret}\}$,
\[
E_{108,a}=\hbar_a\omega_{108},\qquad
m_{108,a}=\frac{E_{108,a}}{c_{\rm int}^2}.
\]
Les longueurs réduite et complète sont respectivement
\[
\frac{\hbar_a}{m_{108,a}c_{\rm int}}=R_{108},
\qquad
\frac{h_a}{m_{108,a}c_{\rm int}}=2\pi R_{108}.
\]
La même action et la même durée interviennent dans les deux lectures
.

\begin{proposition}[Sélecteur de coïncidence des rayons]
Dans la famille dimensionnelle
\[
G_a(\vartheta)=\vartheta\frac{c_{\rm int}^5t_0^2}{\hbar_a},
\qquad \vartheta>0,
\]
la condition de la réalisation circulaire
$G_a(\vartheta)m_{108,a}/c_{\rm int}^2=R_{108}$
admet l'unique solution
\[
\vartheta^\circ_{108}=(54/\pi)^2.
\]
\end{proposition}
\begin{proof}
Le premier rayon est
$r_g(\vartheta)=\vartheta(\pi/54)\ell_0$.
Le comparer à $(54/\pi)\ell_0$, avec $\ell_0>0$, réduit
l'égalité à une équation linéaire strictement croissante.
\end{proof}
La condition de coïncidence est une partie explicite de cette réalisation ;
le théorème résout son coefficient sans utiliser de valeur métrologique de $G$.
Le rayon employé ici est $Gm/c^2$, de sorte que le facteur deux
de la convention $2Gm/c^2$ n'est pas transféré par inadvertance.

\begin{proposition}[Action et gravité réciproques, aire conservée]
Dans les deux coordinatisations précédentes,
\[
\frac{G_{\rm pre}}{G_{\rm ret}}=R_{\rm act}^{-1},\qquad
\hbar_aG_a=\vartheta^\circ_{108}c_{\rm int}^5t_0^2,
\qquad
\ell_P^2=\frac{\hbar_aG_a}{c_{\rm int}^3}
=\vartheta^\circ_{108}\ell_0^2.
\]
\end{proposition}
\begin{proof}
La définition de $G_a$ contient $\hbar_a^{-1}$ et le reste des facteurs
est commun aux deux coordinatisations. Prendre le quotient et le produit donne
les trois égalités.
\end{proof}
Cette réciprocité explique pourquoi l'information d'orientation peut
varier sans que l'aire change. On dispose ainsi d'une relation précise
entre les sections réciproques d'action, l'échelle de Planck et l'opérateur surfacique
de la section~\ref{v:ant:sec:area}.
